% ===================================================================
%  TABEL Nr. 3 — Kinderjaren en jeugd.
%  Traduction 2026 d'apres texto/io, controlee sur texto/fr et
%  texto/es. Consigne : voir texto/nl/10-tabel-01.tex.
% ===================================================================

%%K t03-apar-1 apar
\VUsaut{3.0mm}
\VUtitre{15.0pt}{60}{TABEL Nr. 3}
\VUsaut{1.6mm}
\VUtitre{11.4pt}{0}{Kinderjaren en jeugd.}
\VUsaut{3.4mm}

%%K t03-apar-2 apar
(De 3\textsuperscript{e} en 4\textsuperscript{e} tabellen zijn zo samengesteld dat zij in vier
\VUgras{familietaferelen} de wezenlijke begrippen over het
\VUgras{weer}, de \VUgras{leeftijd}, het \VUgras{gezin}, de \VUgras{kleding} en
de \VUgras{stand} voorstellen.)
\VUsaut{4.0mm}

%%K t03-tit-1 sub
\VUcentre{12.0pt}{0}{\textit{Eerste tafereel.}}
\VUsaut{3.0mm}

%%K t03-tit-2 sub
\VUpk{10.2pt}{40}{Doopplechtigheid in een dorp.}
\VUsaut{2.4mm}

%%K t03-tit-3 sub
\VUpk{10.2pt}{40}{De lente.}
\VUsaut{3.0mm}

%%K t03-c1-01-1 p
1. --- Wij zijn in de \VUgras{lente}, in de maanden \VUgras{maart},
\VUgras{april} of \VUgras{mei}, op de dagen van de \VUgras{week},
\VUgras{maandag}, \VUgras{dinsdag}, of \VUgras{woensdag}. Het is \VUgras{tien
uur} in de morgen.

%%K t03-c1-02-1 p
2. --- Het \VUgras{weer} is mooi, de \VUgras{lucht} zuiver. De zon schijnt. Enkele
\VUgras{wolkjes} \textsuperscript{(2)} stijgen op in de blauwe
\VUgras{hemel} \textsuperscript{(1)}.

%%K t03-c1-03-1 p
3. --- De \VUgras{bomen} hebben nauwelijks \VUgras{bladeren}. De ranke
\VUgras{populieren} \textsuperscript{(3)} en de hoge \VUgras{plataan}
\textsuperscript{(17)} hebben tot nu toe alleen knoppen.

%%K t03-c1-04-1 p
4. --- De \VUgras{zwaluwen} \textsuperscript{(4)} keren terug en fladderen met vrolijk getjilp
om de \VUgras{spits} \textsuperscript{(7)} van de
\VUgras{klokkentoren} \textsuperscript{(5)} die de \VUgras{kerk}
\textsuperscript{(6)} van het dorp bekroont.

%%K t03-tit-4 sub
\VUsaut{3.4mm}
\VUpk{10.2pt}{40}{De kerk.}
\VUsaut{3.0mm}

%%K t03-c1-05-1 p
5. --- Heel eenvoudig is die kerk met haar grote \VUgras{portaal} \textsuperscript{(13)} en
dikke \VUgras{klok} \textsuperscript{(8)}. De kleine
\VUgras{wijzer} \textsuperscript{(9)} staat op het cijfer X, hij wijst de
\VUgras{uren} aan. De \VUgras{grote} \textsuperscript{(10)} staat op XII
(middag); hij wijst de \VUgras{minuten} aan.

%%K t03-c1-06-1 p
6. --- In de klokkentoren luiden de \VUgras{klokken} \textsuperscript{(11)} vrolijk. Op de top
van het dak staat een klein stenen \VUgras{kruis} \textsuperscript{(12)}.

%%K t03-c1-07-1 p
7. --- De kerk heeft niet alleen een groot \VUgras{portaal} \textsuperscript{(13)}, zij heeft
ook een zijdeurtje. Door dat deurtje komt mijnheer \VUgras{pastoor} \textsuperscript{(15)} in
zijn \VUgras{soutane} uit de
\VUgras{sacristie} \textsuperscript{(14)} en gaat hij naar de
\VUgras{pastorie} \textsuperscript{(19)}.

%%K t03-c1-08-1 p
8. --- Naast hem staat de \VUgras{melkvrouw} \textsuperscript{(24)} met haar
\VUgras{ezel} \textsuperscript{(23)}. Zij komt terug uit de stad, waarheen zij
goede melk voor de kinderen bracht.

%%K t03-tit-5 sub
\VUsaut{4.0mm}
\VUpk{10.2pt}{40}{De boomgaard.}
\VUsaut{3.4mm}

%%K t03-c1-09-1 p
9. --- Mijnheer Aliboron (de ezel) is heel tevreden over die ochtendrit. Hij ademt met genot
de lucht in die gebalsemd is door de geur van de
\VUgras{bloemen} \textsuperscript{(20)} die de \VUgras{vruchtbomen}
\textsuperscript{(18)} van de naburige \VUgras{boomgaard} bedekken. Helemaal roze en wit zijn
die \VUgras{perzikbomen} \textsuperscript{(18)} en
\VUgras{abrikozenbomen} \textsuperscript{(19)}, en zij doen een goede oogst
hopen.

%%K t03-c1-10-1 p
10. --- Intussen gaan de kleine \VUgras{bijen} \textsuperscript{(22)} van de
\VUgras{bijenkorf} \textsuperscript{(21)} daarheen om zoemend hun zoete
\VUgras{honing} te roven.

%%K t03-c1-11-1 p
11. --- Maar wat gebeurt er? Daar komen de kinderen van het dorp toelopen en de verschrikte
\VUgras{ganzen} \textsuperscript{(25)} op de vlucht jagen. Dat is het uitgaan van de kerk van
de stoet na het \VUgras{doopsel} van de zoon van de notaris.

%%K t03-c1-12-1 p
12. --- De kleine Jacob wordt in zijn \VUgras{wieg} \textsuperscript{(26)} wakker van het
vrolijke luiden van de klokken, en zijn zusje Magdalena die ook de doopstoet wil zien, vraagt
haar moeder te komen om haar haar te kammen en haar op de drempel van het huis aan te kleden.

%%K t03-c1-13-1 p
13. --- De moeder stemt toe. Zij doet haar \VUgras{hoofddoek} \textsuperscript{(28)}, haar
\VUgras{lijfje} \textsuperscript{(29)} en haar
\VUgras{schort} \textsuperscript{(30)} aan, en zij komt op een klein stoeltje
voor de deur zitten. Magdalena heeft alleen haar \VUgras{onderrok} \textsuperscript{(31)} en
haar \VUgras{korset} \textsuperscript{(32)} aan.

%%K t03-c1-14-1 p
14. --- Wat is zij gelukkig! Zij vergeet haar lievelingsbloemen, haar
\VUgras{anjers} \textsuperscript{(34)}, waarop de \VUgras{vlinders}
\textsuperscript{(35)} neerstrijken, haar \VUgras{tulpen} \textsuperscript{(36)}, haar
\VUgras{lelietjes-van-dalen} \textsuperscript{(37)}, in \VUgras{potten}
\textsuperscript{(38)}, op de \VUgras{muur} \textsuperscript{(33)}, en haar \VUgras{rozen}
\textsuperscript{(39)} die zo'n mooie haag vormen. Zij bekijkt de rijke kleding van de dames
van het gevolg.

%%K t03-c1-15-1 p
15. --- De kinderen van het dorp letten niet zo op de kleding. Zij dringen naar voren en duwen
elkaar om \VUgras{snoepgoed} \textsuperscript{(55)} of
\VUgras{muntstukken} \textsuperscript{(52)} te krijgen. Een van hen huilt
\textsuperscript{(40)}.

%%K t03-c1-16-1 p
16. --- De andere is op de poot van de \VUgras{hond} \textsuperscript{(42)} gestapt die
jankend wegloopt en de \VUgras{hen} \textsuperscript{(43)} en haar
\VUgras{kuikens} \textsuperscript{(44)} op de vlucht jaagt.

%%K t03-tit-6 sub
\VUsaut{5.0mm}
\VUpk{10.2pt}{40}{De doopstoet.}
\VUsaut{4.0mm}

%%K t03-c1-17-1 p
17. --- Voor het gevolg loopt de \VUgras{min} \textsuperscript{(45)}. Zij heeft een mooie
\VUgras{muts} \textsuperscript{(46)} met lange linten. Zij draagt de
\VUgras{pasgeborene} \textsuperscript{(47)} helemaal in \VUgras{kant}
\textsuperscript{(48)} gekleed.

%%K t03-c1-18-1 p
18. --- Achter de min komen de \VUgras{peter} \textsuperscript{(49)} en de
\VUgras{meter} \textsuperscript{(53)}. Zij delen het snoepgoed en de muntstukken
uit. De peter heeft \VUgras{handschoenen} \textsuperscript{(51)} en een
\VUgras{hoge hoed} \textsuperscript{(50)}. De meter houdt een mooi
\VUgras{boeket} \textsuperscript{(54)} in de hand.

%%K t03-c1-19-1 p
19. --- Achter hen zien wij de \VUgras{oom} \textsuperscript{(56)} en de
\VUgras{tante} \textsuperscript{(58)} van het kind. De tante heeft een sierlijke
\VUgras{hoed} \textsuperscript{(59)}, de oom een zwarte \VUgras{geklede jas}
\textsuperscript{(57)} en een wit vest. Daarna komen de \VUgras{neef} \textsuperscript{(60)}
en de \VUgras{nicht} \textsuperscript{(61)}. Deze laatste houdt bij de hand het \VUgras{zusje}
\textsuperscript{(62)} van de pasgeborene, de kleine Bertha.

%%K t03-c1-20-1 p
20. --- Het andere \VUgras{broertje} \textsuperscript{(63)} van Bertha loopt achteraan. Hij
geeft een \VUgras{bloedverwante} \textsuperscript{(64)} de hand. De \VUgras{vrienden}
\textsuperscript{(65)} van het gezin komen daarna.

%%K t03-tit-7 sub
\VUsaut{4.6mm}
\VUpk{10.2pt}{40}{De toeschouwers.}
\VUsaut{3.4mm}

%%K t03-c1-21-1 p
21. --- Naast de stoet staat een \VUgras{bedelaar} \textsuperscript{(66)} die op zijn
\VUgras{kruk} \textsuperscript{(67)} leunt. Hij vraagt een aalmoes.

%%K t03-c1-22-1 p
22. --- Aan de andere kant staat een klein heel \VUgras{beleefd} \textsuperscript{(68)}
jongetje. Hij groet en wenst « \VUgras{goedendag} » aan de mensen die hij kent.

%%K t03-c1-23-1 p
23. --- De dorpelingen zijn hun huizen uitgekomen om de doopstoet te zien. Wij zien een
\VUgras{boer} \textsuperscript{(69)} met zijn \VUgras{katoenen muts} en daarnaast een
\VUgras{boerin} \textsuperscript{(70)}. Daarna een \VUgras{oude vrouw} \textsuperscript{(71)}
die op haar \VUgras{stok} \textsuperscript{(72)} leunt. Ten slotte de \VUgras{molenaar}
\textsuperscript{(73)} met zijn grote
\VUgras{vilten hoed} \textsuperscript{(74)} en een jonge boerin op
\VUgras{klompen} \textsuperscript{(75)}, die haar kindje leert lopen.

%%K t03-c1-24-1 p
24. --- Dat tafereel is heel vermakelijk.
\VUsaut{4.0mm}

%%K t03-tit-8 sub
\VUcentre{12.0pt}{0}{\textit{Tweede tafereel.}}
\VUsaut{3.0mm}

%%K t03-tit-9 sub
\VUpk{10.2pt}{40}{Openbaar feest.}
\VUsaut{2.4mm}

%%K t03-tit-10 sub
\VUpk{10.2pt}{40}{Watersport en roeiwedstrijden.}
\VUsaut{3.0mm}

%%K t03-c2-01-1 p
1. --- De \VUgras{zomer} is gekomen. De hete zon van de maand \VUgras{juni} (juli of
\VUgras{augustus}) spoort de jongelui aan om te baden en te roeien.

%%K t03-c2-02-1 p
2. --- Op de rechteroever van de stroom kleden de roeiers zich uit om aan de watersport en de
roeiwedstrijden deel te nemen.

%%K t03-c2-03-1 p
3. --- Een van hen doet zijn \VUgras{schoenen} \textsuperscript{(1)} uit; hij maakt ze los
(knoopt ze open). Zijn twee kameraden zijn al klaar. Voortaan hebben zij alleen hun
\VUgras{tricot} \textsuperscript{(2)} en
\VUgras{zwembroek} \textsuperscript{(3)} aan. Zij praten terwijl zij op hun
vrienden wachten. Zij spreken over de kermis en het feest dat op de andere oever plaatsvindt.

%%K t03-c2-04-1 p
4. --- Op de grond, naast hen, liggen de \VUgras{kleren} van hun kameraden: een
\VUgras{paar} \VUgras{laarsjes} \textsuperscript{(4)}, \VUgras{schoenen}
\textsuperscript{(5)}, \VUgras{sokken} \textsuperscript{(6)}, \VUgras{vilten}
\textsuperscript{(7)} en \VUgras{strooien} \textsuperscript{(8)} hoeden en een
\VUgras{das} \textsuperscript{(9)}.

%%K t03-c2-05-1 p
5. --- Een van de jongelui heeft zorgvuldig zijn \VUgras{jasje} \textsuperscript{(10)}
opgevouwen. Hij legt het op een handdoek, waarin zijn buurman weldra de twee
\VUgras{kledingstukken} zal wikkelen.

%%K t03-c2-06-1 p
6. --- Een zesde jongen is te laat. Hij heeft nog zijn \VUgras{pet} \textsuperscript{(11)} op
het hoofd. Hij heeft noch zijn \VUgras{hemd} \textsuperscript{(12)}, noch zijn \VUgras{boord}
\textsuperscript{(13)}, noch zijn \VUgras{manchetten} \textsuperscript{(14)} uitgedaan. Hij
houdt zijn
\VUgras{vest} \textsuperscript{(17)} in de hand en heeft nog zijn
\VUgras{broek} \textsuperscript{(16)} en \VUgras{bretels}
\textsuperscript{(15)} aan. Hij zal zeker niet klaar zijn voor de volgende wedstrijd.

%%K t03-c2-07-1 p
7. --- Naast de jongelui leidt een pad met aan weerszijden veel
\VUgras{distels} \textsuperscript{(18)} naar de oever van de rivier, waar vader
Jacob tussen het \VUgras{riet} \textsuperscript{(19)} het \VUgras{vuurwerk}
\textsuperscript{(20)} klaarmaakt dat men 's avonds zal afsteken.

%%K t03-tit-11 sub
\VUsaut{4.4mm}
\VUpk{10.2pt}{40}{De wedstrijd.}
\VUsaut{3.4mm}

%%K t03-c2-08-1 p
8. --- Op de rivier maken enkele dames, die zich met hun parasols beschutten, een tocht in een
\VUgras{boot} \textsuperscript{(21)}. Zij volgen de wedstrijd die heel boeiend is. In elke
\VUgras{jol} \textsuperscript{(22)} roeien vier
\VUgras{roeiers} \textsuperscript{(24)} uit alle macht, terwijl de
\VUgras{stuurman} \textsuperscript{(26)} het \VUgras{roer}
\textsuperscript{(25)} vasthoudt en de broze roeiboot bestuurt. De
\VUgras{riemen} \textsuperscript{(23)} slaan in de maat op het water en de
lichte boten varen met duizelingwekkende snelheid.

%%K t03-c2-09-1 p
9. --- Enkele jongelui wedijveren, naast de pijler van de brug, om de prijs van de
\VUgras{sprietmast} \textsuperscript{(28)}. Een van hen loopt voorzichtig over de buigzame en
gladde mast om het kleine \VUgras{vlaggetje} \textsuperscript{(29)} te bereiken dat aan het
uiteinde bevestigd is. Zijn ongelukkige \VUgras{mededinger} \textsuperscript{(30)} is in het
water gevallen. Hij zwemt om aan land te komen.

%%K t03-c2-10-1 p
10. --- Oudere jongelui houden een \VUgras{steekspel (te water)} \textsuperscript{(32)} naast
de \VUgras{kade} \textsuperscript{(33)}. Al die spelen worden bijgewoond door een talrijk
\VUgras{publiek} \textsuperscript{(31)} dat tegen de \VUgras{leuning} van de \VUgras{brug}
\textsuperscript{(27)} leunt en zich op de \VUgras{oever} verdringt. Enkele toeschouwers keren
zich echter om, om het spel van de \VUgras{prijsmast} \textsuperscript{(45)} te volgen of om
het \VUgras{gemeentebestuur} te bewonderen dat voor de \VUgras{prijsuitreiking} komt.

%%K t03-c2-11-1 p
11. --- Eerst lopen enkele \VUgras{straatjongens} \textsuperscript{(35)} voor de
\VUgras{muzikanten} \textsuperscript{(36)} uit lopen; daarna komen de
\VUgras{burgemeester} \textsuperscript{(37)}, de wethouders en de
\VUgras{gemeenteraad} van het stadje. Ten slotte lopen de
\VUgras{brandweerlieden} \textsuperscript{(38)}.

%%K t03-tit-12 sub
\VUsaut{5.0mm}
\VUpk{10.2pt}{40}{De kermis.}
\VUsaut{3.6mm}

%%K t03-c2-12-1 p
12. --- Talrijke mensen zijn op het \VUgras{kermisplein} \textsuperscript{(39)}. Men komt de
verschillende \VUgras{kramen} zien: de houten \VUgras{paarden} \textsuperscript{(40)}, het
\VUgras{circus} \textsuperscript{(41)} waar de voorstelling al begonnen is; het
\VUgras{openbaar bal} \textsuperscript{(42)}, waar de jongens en de meisjes van de stad het
genoegen van het \VUgras{dansen} smaken.

%%K t03-c2-13-1 p
13. --- Achterin zien wij de \VUgras{arena} \textsuperscript{(44)} waar de dappere Spaanse
\VUgras{stierenvechter} tegen de woedende \VUgras{stier} \textsuperscript{(45)} strijdt,
daarna de \VUgras{rolschaatsbaan} \textsuperscript{(49)}, de \VUgras{schiettent}
\textsuperscript{(46)} waar men zich oefent in het gebruik van de \VUgras{geweren} en in het
schieten op de
\VUgras{schijf}.

%%K t03-c2-14-1 p
14. --- Vlakbij staat het \VUgras{kermistheater} \textsuperscript{(47)} waarin men het
\VUgras{blijspel} en het \VUgras{drama} speelt. Vlakbij staat de kraam van de \VUgras{reus}
\textsuperscript{(48)} en van de \VUgras{dwerg}. De
\VUgras{lengte} van de eerste is twee meter vijftien centimeter; de tweede is
nauwelijks zo groot als een kind.

%%K t03-c2-15-1 p
15. --- En helemaal achteraan staat de grote \VUgras{menagerie} \textsuperscript{(50)} met
haar \VUgras{wilde dieren}, haar \VUgras{tijger} \textsuperscript{(51)} met krachtige
\VUgras{klauwen}, haar \VUgras{leeuw} \textsuperscript{(52)} met dichte \VUgras{manen}, haar
twee \VUgras{giraffen} \textsuperscript{(53)} en haar \VUgras{olifant} \textsuperscript{(54)}
die zo behendig zijn
\VUgras{slurf} gebruikt.

%%K t03-c2-16-1 p
16. --- De \VUgras{temmer} \textsuperscript{(55)} die die dieren africht staat bij de deur van
de kraam.
\VUsaut{6.0mm}
\VUfilet{22mm}
